\subsection{Bloom filters library for large batches of items}
\label{sec:bfs}

\paragraph{Presentation}
BFlib is a Bloom filters library aimed at storing and querying large chunks of elements at once.
Instead of one big array, the library uses a collection of smaller filters.
The library is specifically designed to run with the Processing-in-Memory hardware from UPMEM, in which case the many filters actually reside on different memory banks.
Efficient parallel processing can thus be performed to execute large batches of insertions or lookups.
We re-implement this processing using the BPL library.

\paragraph{Limitations encountered} Since we execute insertions and lookups in batches, the filters need to exist in memory across several executions.
However, the BPL does not support such feature at the moment.
Therefore, we had to design an additional class to manage persistent storage and implement it for all architectures supported by the BPL.
For UPMEM, the persistent data block is reserved in the MRAM of each DPU.
As for Multicore, we allocate all data blocks required, store them in a vector, and adds an index argument to the task so that each instance can retrieve its assigned block.

\paragraph{Benchmarks} We consider Bloom filters using 8 hash functions, and a varying size (128 KiB or 1 MiB).
For PiM implementations, we allocate 8 ranks, i.e. 512 processing units.
One benchmark consists of inserting a varying number of items (10M or 100M), and then querying all the inserted items.

\begin{table}[htbp]
    \centering
    \caption{Results of the Bloom filters benchmarks}
    \label{tab:bflib-results}
    \begin{tabular}{lrrrrrr}
    \toprule
     & & & \multicolumn{2}{c}{Multicore time (s)} & \multicolumn{2}{c}{PiM time (s)} \\
    \cmidrule(l{3pt}r{3pt}){4-5}\cmidrule(l{3pt}r{3pt}){6-7}
    & Filter size & \#items & Standard* & BPL & SDK & BPL \\
    \midrule
    \multirow{4}{*}{Insertions} & 128 MiB & 10M & 0.22 & 0.56 & 0.13 & 0.11 \\
    & 128 MiB & 100M & 1.49 & 4.64 & 0.66 & 0.68 \\
    & 1 GiB & 10M & 0.39 & 0.84 & 0.13 & 0.11 \\
    & 1 GiB & 100M & 2.24 & 5.56 & 0.66 & 0.68 \\
    \midrule
    \multirow{4}{*}{Lookups} & 128 MiB & 10M & 0.25 & 0.41 & 0.27 & 0.28 \\
    & 128 MiB & 100M & 2.26 & 3.06 & 2.13 & 2.26 \\
    & 1 GiB & 10M & 0.38 & 0.38 & 0.27 & 0.28 \\
    & 1 GiB & 100M & 2.89 & 2.94 & 2.13 & 2.81 \\
    \bottomrule
    \end{tabular}
    \\ \vspace{0.3em}
    {\scriptsize *Uses a single big filter and not the many small filters approach}
\end{table}

\paragraph{Results} We present the results in \tablename~\ref{tab:bflib-results}.
The Standard Multicore implementation appears faster than the BPL Multicore version, due to being less memory-bound.
Indeed, it handles items one at a time instead of regrouping them in batches and processing them altogether.
Batching is beneficial with the PiM architecture to minimize the overhead of transferring data between the main memory and the small memory banks, but this purpose is absent on Multicore.
Implementing a program on several architectures at once is a great perk of the BPL, but this comes with a caveat: it forces to use the same parallelization scheme for both Multicore and PiM, which is not necessarily the best choice performance-wise.
Regarding PiM, the BPL version has similar execution times to the SDK program.

\paragraph{Source code} The BPL implementation is publicly available at \url{https://gitlab.inria.fr/pim/org.pim.bflib-bpl} (GNU AGPLv3).
The original Standard Multicore and SDK implementations are available at \url{https://gitlab.inria.fr/pim/org.pim.bloomfilters} (GNU AGPLv3).